% GENERATED FILE. Edit canonical lessons, then run npm run generate:books.
% canonical-source-hash: fnv1a64:aaad1bcf94a77174
% canonical-lessons: ML-C522-jijnasa, ML-C522-aksama, ML-C522-desyappetuka, ML-C522-unmesam, ML-C522-matukkuka

\chapter{Curious and Fed Up}
\label{ch:ml-curious-and-fed-up}
\hlchaptermodality{522}

\begin{chapteropening}
\textbf{By the end of this chapter:} \emph{I can say curiosity, impatience, to get angry, freshness and to get fed up.}

Complete the last lesson of chapter 522: I can say curiosity, impatience, to get angry, freshness and to get fed up.
\end{chapteropening}

\section[\texorpdfstring{\ml{ജിജ്ഞാസ} (jijñāsa)}{ജിജ്ഞാസ (jijñāsa)}]{\ml{ജിജ്ഞാസ} (jijñāsa) — curiosity}

\label{lesson:ML-C522-jijnasa}



\begin{quote}
Before the new one: say the Malayalam for nostalgia, then the Malayalam for to sigh.
\end{quote}

\subsection*{You'll want to know: \ml{ജിജ്ഞാസ}}
\textbf{\ml{ജിജ്ഞാസ}} — \emph{jijñāsa} — \textquotedblleft{}curiosity\textquotedblright{}.

\begin{quote}
\ml{കുട്ടികൾക്ക്} \ml{എല്ലാ} \ml{കാര്യത്തിലും} \ml{ജിജ്ഞാസയുണ്ട്}.
\end{quote}

\begin{quote}
\emph{kuṭṭikaḷkkŭ ellā kāryattiluṁ jijñāsayuṇṭŭ.}
\end{quote}

\begin{quote}
Children are curious about everything.
\end{quote}

What one is curious about takes -\ml{ഇൽ}. \ml{ജിജ്ഞാസയുള്ള}, curious, is the adjective.

\begin{cousinweb}
\textbf{\ml{ജിജ്ഞാസ}} (\emph{jijñāsa}) is borrowed from Sanskrit \emph{jijñāsā}, \textquotedblleft{}the wish to know\textquotedblright{} (Monier-Williams), the wish-form of the root \emph{jñā}, \textquotedblleft{}to know\textquotedblright{}.
\end{cousinweb}

\begin{grammarlens}[title={What you've built}]
One more word for describing what you lived and felt.
\end{grammarlens}

\subsection*{Your turn}
\begin{itemize}
\raggedright
  \item \emph{Say it:} \emph{jijñāsa}
  \item \emph{Say it:} \emph{jijñāsa}, once more
  \item \emph{Say it:} \emph{kuṭṭikaḷkkŭ ellā kāryattiluṁ jijñāsayuṇṭŭ.} — Children are curious about everything.
  \item \emph{Recall:} say the Malayalam for nostalgia, then the Malayalam for to sigh, then say \emph{jijñāsa} again
\end{itemize}

\begin{tcolorbox}[breakable,colback=teal!4,colframe=teal!35!black,title={Before you move on}]
What does \ml{ജിജ്ഞാസ} mean? (\textquotedblleft{}Curiosity\textquotedblright{}.) Say it once more.
\end{tcolorbox}
\section[\texorpdfstring{\ml{അക്ഷമ} (akṣama)}{അക്ഷമ (akṣama)}]{\ml{അക്ഷമ} (akṣama) — impatience}

\label{lesson:ML-C522-aksama}



\begin{quote}
Before the new one: say the Malayalam for to sigh, then the Malayalam for curiosity.
\end{quote}

\subsection*{You'll want to know: \ml{അക്ഷമ}}
\textbf{\ml{അക്ഷമ}} — \emph{akṣama} — \textquotedblleft{}impatience\textquotedblright{}.

\begin{quote}
\ml{ബസ്} \ml{വൈകിയപ്പോൾ} \ml{യാത്രക്കാരുടെ} \ml{അക്ഷമ} \ml{കൂടി}.
\end{quote}

\begin{quote}
\emph{basŭ vaikiyappōḷ yātrakkāruṭe akṣama kūṭi.}
\end{quote}

\begin{quote}
As the bus got later, the passengers' impatience grew.
\end{quote}

\ml{അ}- not + \ml{ക്ഷമ} patience, the word inside \ml{ക്ഷമിക്കണം}.

\begin{cousinweb}
\textbf{\ml{അക്ഷമ}} (\emph{akṣama}) is borrowed from Sanskrit \emph{akṣamā}, \textquotedblleft{}impatience\textquotedblright{} (Monier-Williams), \emph{a-}, \textquotedblleft{}not\textquotedblright{}, and \emph{kṣamā}, \textquotedblleft{}patience\textquotedblright{}.
\end{cousinweb}

\begin{grammarlens}[title={What you've built}]
One more word for describing what you lived and felt.
\end{grammarlens}

\subsection*{Your turn}
\begin{itemize}
\raggedright
  \item \emph{Say it:} \emph{akṣama}
  \item \emph{Say it:} \emph{akṣama}, once more
  \item \emph{Say it:} \emph{basŭ vaikiyappōḷ yātrakkāruṭe akṣama kūṭi.} — As the bus got later, the passengers' impatience grew.
  \item \emph{Recall:} say the Malayalam for to sigh, then the Malayalam for curiosity, then say \emph{akṣama} again
\end{itemize}

\begin{tcolorbox}[breakable,colback=teal!4,colframe=teal!35!black,title={Before you move on}]
What does \ml{അക്ഷമ} mean? (\textquotedblleft{}Impatience\textquotedblright{}.) Say it once more.
\end{tcolorbox}
\section[\texorpdfstring{\ml{ദേഷ്യപ്പെടുക} (dēṣyappeṭuka)}{ദേഷ്യപ്പെടുക (dēṣyappeṭuka)}]{\ml{ദേഷ്യപ്പെടുക} (dēṣyappeṭuka) — to get angry; to lose one's temper}

\label{lesson:ML-C522-desyappetuka}



\begin{quote}
Before the new one: say the Malayalam for curiosity, then the Malayalam for impatience.
\end{quote}

\subsection*{You'll want to know: \ml{ദേഷ്യപ്പെടുക}}
\textbf{\ml{ദേഷ്യപ്പെടുക}} — \emph{dēṣyappeṭuka} — \textquotedblleft{}to get angry; to lose one's temper\textquotedblright{}.

\begin{quote}
\ml{ചെറിയ} \ml{കാര്യത്തിന്} \ml{അവൻ} \ml{ദേഷ്യപ്പെട്ടു}.
\end{quote}

\begin{quote}
\emph{ceṟiya kāryattinŭ avan dēṣyappeṭṭu.}
\end{quote}

\begin{quote}
He got angry over a small matter.
\end{quote}

Past: \ml{ദേഷ്യപ്പെട്ടു}, \emph{dēṣyappeṭṭu}.

\ml{ദേഷ്യം} anger + \ml{പെടുക}, the verb that turns a feeling into something that happens to you. The person you are angry with takes -\ml{ഓട്}.

\begin{cousinweb}
\textbf{\ml{ദേഷ്യപ്പെടുക}} (\emph{dēṣyappeṭuka}) is built from \ml{ദേഷ്യം}, \textquotedblleft{}anger\textquotedblright{}, which comes from Sanskrit \emph{dveṣa}, \textquotedblleft{}hatred\textquotedblright{} (CDIAL 6696), and the native helper verb \ml{പെടുക}.
\end{cousinweb}

\begin{grammarlens}[title={What you've built}]
One more word for describing what you lived and felt.
\end{grammarlens}

\subsection*{Your turn}
\begin{itemize}
\raggedright
  \item \emph{Say it:} \emph{dēṣyappeṭuka}
  \item \emph{Say it:} \emph{dēṣyappeṭuka}, once more
  \item \emph{Say it:} \emph{ceṟiya kāryattinŭ avan dēṣyappeṭṭu.} — He got angry over a small matter.
  \item \emph{Recall:} say the Malayalam for curiosity, then the Malayalam for impatience, then say \emph{dēṣyappeṭuka} again
\end{itemize}

\begin{tcolorbox}[breakable,colback=teal!4,colframe=teal!35!black,title={Before you move on}]
What does \ml{ദേഷ്യപ്പെടുക} mean? (\textquotedblleft{}To get angry; to lose one's temper\textquotedblright{}.) Say it once more.
\end{tcolorbox}
\section[\texorpdfstring{\ml{ഉന്മേഷം} (unmēṣaṁ)}{ഉന്മേഷം (unmēṣaṁ)}]{\ml{ഉന്മേഷം} (unmēṣaṁ) — freshness; vigour}

\label{lesson:ML-C522-unmesam}



\begin{quote}
Before the new one: say the Malayalam for impatience, then the Malayalam for to get angry.
\end{quote}

\subsection*{You'll want to know: \ml{ഉന്മേഷം}}
\textbf{\ml{ഉന്മേഷം}} — \emph{unmēṣaṁ} — \textquotedblleft{}freshness; vigour\textquotedblright{}.

\begin{quote}
\ml{രാവിലെ} \ml{ഒരു} \ml{ചായ} \ml{കുടിച്ചാൽ} \ml{എനിക്ക്} \ml{ഉന്മേഷം} \ml{കിട്ടും}.
\end{quote}

\begin{quote}
\emph{rāvile oru cāya kuṭiccāl enikkŭ unmēṣaṁ kiṭṭuṁ.}
\end{quote}

\begin{quote}
If I drink a tea in the morning, I feel fresh.
\end{quote}

The fresh, wide-awake feeling after rest or a bath: \ml{കുളി} \ml{കഴിഞ്ഞ്} \ml{ഉന്മേഷം} \ml{തോന്നി}.

\begin{cousinweb}
\textbf{\ml{ഉന്മേഷം}} (\emph{unmēṣaṁ}) is borrowed from Sanskrit \emph{unmeṣa}, \textquotedblleft{}the opening of the eyes\textquotedblright{} (Monier-Williams). Malayalam took it for the feeling of being wide awake and fresh.
\end{cousinweb}

\begin{grammarlens}[title={What you've built}]
One more word for describing what you lived and felt.
\end{grammarlens}

\subsection*{Your turn}
\begin{itemize}
\raggedright
  \item \emph{Say it:} \emph{unmēṣaṁ}
  \item \emph{Say it:} \emph{unmēṣaṁ}, once more
  \item \emph{Say it:} \emph{rāvile oru cāya kuṭiccāl enikkŭ unmēṣaṁ kiṭṭuṁ.} — If I drink a tea in the morning, I feel fresh.
  \item \emph{Recall:} say the Malayalam for impatience, then the Malayalam for to get angry, then say \emph{unmēṣaṁ} again
\end{itemize}

\begin{tcolorbox}[breakable,colback=teal!4,colframe=teal!35!black,title={Before you move on}]
What does \ml{ഉന്മേഷം} mean? (\textquotedblleft{}Freshness; vigour\textquotedblright{}.) Say it once more.
\end{tcolorbox}
\section[\texorpdfstring{\ml{മടുക്കുക} (maṭukkuka)}{മടുക്കുക (maṭukkuka)}]{\ml{മടുക്കുക} (maṭukkuka) — to get fed up; to tire of}

\label{lesson:ML-C522-matukkuka}



\begin{quote}
Before the new one: say the Malayalam for to get angry, then the Malayalam for freshness.
\end{quote}

\subsection*{You'll want to know: \ml{മടുക്കുക}}
\textbf{\ml{മടുക്കുക}} — \emph{maṭukkuka} — \textquotedblleft{}to get fed up; to tire of\textquotedblright{}.

\begin{quote}
\ml{ഒരേ} \ml{ഭക്ഷണം} \ml{കഴിച്ച്} \ml{എനിക്ക്} \ml{മടുത്തു}.
\end{quote}

\begin{quote}
\emph{orē bhakṣaṇaṁ kaḻiccŭ enikkŭ maṭuttu.}
\end{quote}

\begin{quote}
I am fed up with eating the same food.
\end{quote}

Past: \ml{മടുത്തു}, \emph{maṭuttu}.

The one who is fed up takes -\ml{ക്ക്}: \ml{എനിക്ക്} \ml{മടുത്തു}. The past is \ml{മടുത്തു}.

\begin{cousinweb}
\textbf{\ml{മടുക്കുക}} (\emph{maṭukkuka}) is native Dravidian: DEDR 4652 lists the Malayalam verb itself, to be tired of or to loathe, beside \ml{മടിക്കുക}, \textquotedblleft{}to grow lazy, to hold back\textquotedblright{} (DEDR 4652).
\end{cousinweb}

\begin{grammarlens}[title={What you've built}]
One more word for describing what you lived and felt.
\end{grammarlens}

\subsection*{Your turn}
\begin{itemize}
\raggedright
  \item \emph{Say it:} \emph{maṭukkuka}
  \item \emph{Say it:} \emph{maṭukkuka}, once more
  \item \emph{Say it:} \emph{orē bhakṣaṇaṁ kaḻiccŭ enikkŭ maṭuttu.} — I am fed up with eating the same food.
  \item \emph{Recall:} say the Malayalam for to get angry, then the Malayalam for freshness, then say \emph{maṭukkuka} again
\end{itemize}

\begin{tcolorbox}[breakable,colback=teal!4,colframe=teal!35!black,title={Before you move on}]
What does \ml{മടുക്കുക} mean? (\textquotedblleft{}To get fed up; to tire of\textquotedblright{}.) Say it once more.
\end{tcolorbox}
